\section[Structural and Signed-Center Optimization]
{Structural, Shared-Local, and Signed-Center Optimization}
\label{app:structural-shared-local-signed-center}

This appendix solves the coordinate and compactness problems used by the
geometric interfaces in Section~\ref{sec:structure-common}.  It also proves
the shared equilateral-enclosure calculus used by the later optimization
appendices.  Each calculation is attached to the body statement it supplies.

\subsection{Open--closed shrinking feasibility}

\begin{lemma}[Uniform shrinking margin]
\label{lem:app-open-closed-shrinking}
If seven open unit equilateral triangles cover $H$, then, for some
$\varepsilon>0$, seven closed equilateral triangles of side
$1-\varepsilon$ also cover $H$.
\end{lemma}

\begin{proof}
For $p\in H$ put
\[
 m_\nu(p):=\operatorname{dist}(p,\mathbb R^2\setminus U_\nu)
 \quad\bigl(\nu\in\{C,0,\ldots,5\}\bigr),
 \qquad
 m(p):=\max_{\nu\in\{C,0,\ldots,5\}}m_\nu(p).
\]
The continuous function $m$ is positive on compact $H$, hence
\[
 \eta:=\min_{p\in H}m(p)>0.
\]
Choose
\[
 0<\rho<\min\left\{\eta,\frac{\sqrt3}{6}\right\}.
\]
Moving every side of $U_\nu$ inward by $\rho$ gives a closed equilateral
triangle $K_\nu$ of side $1-2\sqrt3\rho$.  For each $p\in H$, some
$m_\nu(p)\ge\eta>\rho$, so $p\in K_\nu$.  Thus
$\varepsilon=2\sqrt3\rho$ is feasible.  This supplies the only quantitative
step in Proposition~\ref{prop:open-closed-scaled}.
\end{proof}

\subsection{T3-like trace-majorant feasibility}

The calculation uses the local wedge and the two orientation charts developed
in the exact-trace module below.  It is placed here beside the body warning
that the output is a closed majorant, not a replacement open triangle.

\begin{lemma}[T3-like translation calculation]
\label{lem:app-t3-translation-calculation}
The closed-trace majorant in Proposition~\ref{prop:t3-translation} exists.
In the calculation chart, every original $\Tlike$ triangle also satisfies
$A_i+B_i\le1$.
\end{lemma}

\begin{proof}
Use the wedge chart of \zcref{lem:app-local-wedge-calculation} and the
orientation forms \eqref{eq:orientation-type-I}--
\eqref{eq:orientation-type-II} in
Subsection~\ref{subsec:app-exact-trace-normalization}.  In Type I,
direct inversion shows that when exactly two vertices lie outside $W$, the
sole wedge vertex has both coordinates strictly below $1$: the relevant
bounds are $(1-t)/z$ and $t/z$, with strictness from $\alpha>0$ or
$\beta>0$.  The support argument in
\zcref{lem:app-local-wedge-calculation} then excludes positive-length contact
with either adjacent radial line.  Thus a $\Tlike$ triangle has Type II.

Reflect in $x=y$ if necessary so that $T_0\cap r_1$ has positive length.
The endpoint cases $t=0,1$ give no such interval, hence $0<t<1$.  The exact
axis reaches are
\[
 A_0=\beta,\qquad B_0=z-\alpha-\beta,
\]
and the support condition is
\[
 Q_x>1,\qquad Q_y\le1,
\]
where $Q$ is \eqref{eq:orientation-type-II-Q}.  In particular
\[
 A_0+B_0=z-\alpha<z\le1.
\]

Define $\widehat T$ by replacing $\alpha$ with $0$ and retaining $t,\beta$.
This changes only the translation.  At the origin,
$\widehat U=0$ and $0<\widehat V=\beta<z$, so $V_0$ lies in the relative
interior of a side.  For $(x,y)\in W$,
\[
 \widehat U=(1-t)x+ty\ge0,
 \qquad\widehat V=V,
 \qquad\widehat U+\widehat V=U+V-\alpha.
\]
Thus $T\cap W\subset\widehat T\cap W$.  The old $x$-axis reach is $B_0$ and
the new reach is $B_0+\alpha$, so the inclusion is strict in $H$.

The two outside vertices remain outside, while $V_0$ and the new wedge vertex
lie in $H$.  Moreover $\widehat Q_x>Q_x>1$, and $Q_x>1$ gives
\[
 tD_t(1-\widehat Q_y)
 >D_t(1-z)+(1-t)\alpha>0.
\]
Hence $\widehat Q_y<1$ and $(o(\widehat T),n(\widehat T))=(2,1)$.  At every
nonzero point of $W$, $\widehat U>0$; the other two displayed slack relations
show that old interior points other than $V_0$ remain interior.  Reflection
proves the other supported-arm case.
\end{proof}

\subsection{Strict-handoff feasibility}

\begin{lemma}[Cyclic strict-handoff selection]
\label{lem:app-strict-handoff-selection}
Under the hypotheses of Proposition~\ref{prop:strict-handoffs}, all six
overlap intervals are nonempty, and the selectors can preserve the stated
one- or two-supercritical-index alternatives.
\end{lemma}

\begin{proof}
Since $V_i$ is interior to a diameter-one triangle,
$0<A_i<1$ and $0<B_i<1$.  No nonincident V triangle contains a
relative-interior point of $e_{i,i+1}$.  Thus the relatively open incident
traces cover the edge and overlap, so
\[
 (1-A_{i+1},B_i)\ne\varnothing.
\]
Choosing $\xi_i$ in this interval places both selected anchors in $U_i$.
Their squared distance is $a_i^2+a_ib_i+b_i^2$, which is strictly below $1$
because both are interior points of the same open unit triangle.

Put $\ell_i:=1-A_{i+1}$ and $u_i:=B_i$.  Then
\begin{equation}
 A_i+B_i>1\iff\ell_{i-1}<u_i,
 \qquad a_i+b_i>1\iff\xi_{i-1}<\xi_i.
 \label{eq:actual-selected-ascents}
\end{equation}
If $i$ is actually nonsupercritical, then
\[
 \xi_i<u_i\le\ell_{i-1}<\xi_{i-1},
\]
so no selection makes it supercritical.  Also
\begin{equation}
 \sum_{i=0}^5(a_i+b_i)
 =\sum_{i=0}^5(1-\xi_{i-1}+\xi_i)=6.
 \label{eq:handoff-telescoping}
\end{equation}
If $p$ is the sole actual supercritical index, the other five selected sums
are below $1$, and \eqref{eq:handoff-telescoping} forces the $p$th sum above
$1$.

For two nonadjacent supercritical indices $p,q$, choose
\[
 \ell_{r-1}<z_r<u_r\qquad(r\in\{p,q\}),
\]
and then
\[
 \xi_{r-1}\in(\ell_{r-1},\min\{u_{r-1},z_r\}),
 \qquad
 \xi_r\in(\max\{\ell_r,z_r\},u_r).
\]
The pairs are disjoint and create ascents at $p,q$.  Any three indices on a
six-cycle contain two nonadjacent indices.  If the only two supercritical
indices are adjacent, say $p,p+1$, the remaining inequalities give
\[
 \ell_{p-1}<\ell_{p+1},
 \qquad
 \max\{\ell_{p-1},\ell_p\}<\min\{u_p,u_{p+1}\}.
\]
Choose $\xi_p$ between the latter bounds and choose
$\xi_{p-1}<\xi_p<\xi_{p+1}$ in their overlap intervals.  By
\eqref{eq:actual-selected-ascents}, both selected indices remain
supercritical.
\end{proof}

\subsection{Shared equilateral-enclosure optimization}
\label{sec:universal-calculus}

Let $\mathsf R$ denote counterclockwise rotation through $2\pi/3$.  For a
nonempty compact $K\subset\mathbb R^2$, put
\[
 h_K(n):=\max_{x\in K}\langle x,n\rangle.
\]

\begin{proposition}[Equilateral enclosure gauge]
\label{prop:new-enclosure-gauge}
\label{prop:universal-enclosure-gauge}
The least side length $\Lambda(K)$ of a closed equilateral triangle containing
$K$ is
\begin{equation}
 \Lambda(K)=\frac2{\sqrt3}
 \min_{\lVert n\rVert=1}\sum_{j=0}^2h_K(\mathsf R^jn).
 \label{eq:universal-lambda}
\end{equation}
It is monotone under inclusion, translation invariant, and positively
homogeneous.  If $K$ is a polygonal hull with at least two distinct points, a
minimizing normal can be chosen perpendicular to a hull edge.
\end{proposition}

\begin{proof}
For fixed outward normals $n,\mathsf Rn,\mathsf R^2n$, the least containing
triangle has the corresponding three support numbers, and its side is
$2/\sqrt3$ times their sum.  Translation invariance follows from
\[
 n+\mathsf Rn+\mathsf R^2n=0,
\]
while monotonicity and homogeneity are immediate.  On an open angular cell
with fixed support points, the support sum is
$A\cos\theta+B\sin\theta$ and its second derivative is its negative.  For a
nonsingleton the support sum is positive, so an interior critical point is a
maximum; a minimum occurs at a support tie, hence at a hull-edge normal.
\end{proof}

\subsubsection{Finite caliper criterion}
\label{subsec:finite-caliper}

Let $\mathsf J$ denote counterclockwise rotation through $\pi/2$.  For a
nonempty finite set $F$ define
\begin{equation}
 W_F(N):=\sum_{k=0}^2\max_{z\in F}
 \langle z,\mathsf R^kN\rangle.
 \label{eq:caliper-support-sum}
\end{equation}

\begin{theorem}[Finite caliper certificate]
\label{thm:cert-caliper}
If $F$ has at least two distinct points, then $F$ is contained in a closed
unit equilateral triangle if and only if there are distinct $p,q\in F$ and
$\varepsilon\in\{-1,1\}$ such that
\begin{equation}
 W_F\bigl(\varepsilon\mathsf J(q-p)\bigr)
 \le\frac{\sqrt3}{2}\lVert p-q\rVert.
 \label{eq:finite-caliper-test}
\end{equation}
A singleton is contained in equilateral triangles of arbitrarily small
positive side length.  For $F_x=\{O,A,B,x\}$, every clause is a finite
polynomial clause of degree at most four in the coordinates of $A,B,x$; at
most $12\cdot4^3$ clauses are needed.
\end{theorem}

\begin{proof}
Apply Proposition~\ref{prop:universal-enclosure-gauge} to $\conv F$.  A
minimizing normal is perpendicular to a hull edge $[p,q]$ and, after cyclic
rotation, equals
$\varepsilon\mathsf J(q-p)/\lVert p-q\rVert$.  Homogeneity gives
\eqref{eq:finite-caliper-test}; a nondegenerate segment follows by checking
its two normals.

For the polynomial assertion fix $(p,q,\varepsilon)$, put
$N=\varepsilon\mathsf J(q-p)$ and $N_k=\mathsf R^kN$, and select support
labels $s_k\in F_x$.  Retain the strict condition
\[
 \lVert p-q\rVert^2>0.
\]
On the cell
\begin{equation}
 \langle s_k-y,N_k\rangle\ge0
 \qquad(y\in F_x,\ k=0,1,2),
 \label{eq:caliper-cell}
\end{equation}
the test becomes
\begin{equation}
 \sum_{k=0}^2\langle s_k,N_k\rangle
 \le\frac{\sqrt3}{2}\lVert p-q\rVert.
 \label{eq:caliper-cell-polynomial}
\end{equation}
The left side is polynomial of degree at most two.  Both sides are
nonnegative on an active cell, so one squaring gives degree at most four.
There are at most six unordered pairs, two normal signs for each pair, and
$4^3$ support triples for each pair-normal candidate.
\end{proof}

\subsection{Center and V triangle incidence feasibility}

\begin{lemma}[Separation of nonadjacent boundary edges]
\label{lem:app-center-edge-separation}
Relative-interior points on two nonadjacent edges of $H$ are more than unit
distance apart.
\end{lemma}

\begin{proof}
Up to symmetry take
\[
 x=(1-t)V_0+tV_1,\qquad 0<t<1.
\]
For $y=(1-s)V_2+sV_3$, $0<s<1$, direct expansion gives
\[
 \lVert x-y\rVert^2-1=(1-t)(2-t)+s(s+t)>0.
\]
For the other orbit, $y=(1-s)V_3+sV_4$, put $u=t+s$ and obtain
\[
 \lVert x-y\rVert^2-1=(u-1)^2+2>0.
\]
These are the two nonadjacent-edge configurations.  Since the diameter of a
unit equilateral triangle is $1$, the result supplies
Proposition~\ref{prop:ce-classification}.
\end{proof}

\begin{lemma}[Local wedge calculation]
\label{lem:app-local-wedge-calculation}
The incidence conclusions of Lemma~\ref{lem:local-wedge} hold.
\end{lemma}

\begin{proof}
Normalize to $i=0$ and use the affine chart
\[
 X=V_0+x(V_1-V_0)+y(V_5-V_0).
\]
Then $V_0=(0,0)$,
\[
 \lVert(x,y)\rVert^2=x^2+y^2-xy,
\]
and, in the closed unit ball about $V_0$, membership in $H$ is membership in
\[
 W:=\{(x,y):x\ge0,\ y\ge0\}.
\]
Indeed,
\[
 x^2+y^2-xy=\left(y-\frac x2\right)^2+\frac{3x^2}{4}
             =(x-y)^2+xy.
\]
The adjacent radial segments are
\[
 r_1=\{(1,y):0\le y\le1\},\qquad
 r_5=\{(x,1):0\le x\le1\}.
\]
Suppose $T_0\cap r_1$ has positive length and let $m$ be the largest
$x$-coordinate of a vertex of $T_0$.  If $m>1$, a maximizing vertex $(m,y)$
lies in the unit ball and has $0<y<1$, so it lies in $H$.  If $m=1$,
positive-length contact with the support line $x=1$ forces a side there, and
both endpoints satisfy $1+y^2-y\le1$ and lie in $H$.  The argument for
$r_5$ is symmetric.  If both arms are met but only one vertex lies in $H$,
that vertex would have $x>1$ and $y>1$, contradicting
$(x-y)^2+xy>1$.

Finally, if all three vertices lay in convex $H$, then $T_0\subset H$,
contrary to the extreme point $V_0$ lying in $\interior(T_0)$.  Hence
$o(T_0)\ge1$, completing the calculation.
\end{proof}

\subsection{Exact-trace normalization}
\label{subsec:app-exact-trace-normalization}

Set
\[
 D_t:=1-t+t^2,\qquad z:=\sqrt{D_t},\qquad0\le t\le1.
\]
After relabelling the vertices, every orientation belongs to one of the two
families
\[
\begin{array}{c|cc}
 &\text{first side vector}&\text{second side vector}\\ \hline
\mathrm I&z^{-1}(1,t)&z^{-1}(1-t,1)\\
\mathrm {II}&z^{-1}(t,-(1-t))&z^{-1}(1,t).
\end{array}
\]
In Type I write
\begin{equation}
 U=\alpha+y-tx,\qquad V=\beta+x-(1-t)y,\qquad
 T=\{U\ge0,V\ge0,U+V\le z\},
 \label{eq:orientation-type-I}
\end{equation}
and in Type II write
\begin{equation}
 U=\alpha+(1-t)x+ty,\qquad V=\beta+tx-y,\qquad
 T=\{U\ge0,V\ge0,U+V\le z\}.
 \label{eq:orientation-type-II}
\end{equation}
In both charts
\begin{equation}
 V_0\in\interior(T)\iff
 \alpha>0,\quad\beta>0,\quad\alpha+\beta<z.
 \label{eq:orientation-interior}
\end{equation}
The Type-II vertex in the wedge is
\begin{equation}
 Q=\left(
 \frac{z-\alpha-t\beta}{D_t},
 \frac{t(z-\alpha)+(1-t)\beta}{D_t}
 \right).
 \label{eq:orientation-type-II-Q}
\end{equation}

\begin{lemma}[Raw $(3,0)$ trace-normalization calculation]
\label{lem:app-vd0-trace-normalization-calculation}
The translation asserted in
Proposition~\ref{prop:vd0-exact-trace-normalization} exists and preserves both
the open and closed traces in $H$.
\end{lemma}

\begin{proof}
Normalize to $i=0$.  A Type-II triangle cannot have $o=3$, because the point
$Q$ in \eqref{eq:orientation-type-II-Q} lies in the wedge; together with the
diameter-one condition, Lemma~\ref{lem:local-wedge} puts it in $H$.  Hence
$T$ has Type-I form.  Its vertices are
\[
 Q_0=\left(-\frac{(1-t)\alpha+\beta}{D_t},
           -\frac{\alpha+t\beta}{D_t}\right),
\]
\begin{equation}
\begin{aligned}
 Q_1&=\left(\frac{z-(1-t)\alpha-\beta}{D_t},
             \frac{tz-\alpha-t\beta}{D_t}\right),\\
 Q_2&=\left(\frac{(1-t)z-(1-t)\alpha-\beta}{D_t},
             \frac{z-\alpha-t\beta}{D_t}\right).
\end{aligned}
\label{eq:o3-type-I-vertices}
\end{equation}
Thus $o=3$ is equivalent to
\begin{equation}
 \alpha+t\beta>tz,\qquad
 (1-t)\alpha+\beta>(1-t)z,
 \label{eq:o3-outside-conditions}
\end{equation}
which forces $0<t<1$.  Put $s=\alpha+\beta$ and $L=z-s>0$.  Then
\begin{equation}
 \alpha>\frac{tL}{1-t},\qquad
 \beta>\frac{(1-t)L}{t}.
 \label{eq:o3-redundancy-conditions}
\end{equation}
For $(x,y)\in W$ satisfying
\begin{equation}
 (1-t)x+ty\le L,
 \label{eq:o3-wedge-cut}
\end{equation}
the minimum of $U$ is attained at $(L/(1-t),0)$ and that of $V$ at
$(0,L/t)$.  Hence
\begin{equation}
 T\cap W=\{(x,y)\in W:(1-t)x+ty\le L\}.
 \label{eq:o3-wedge-trace}
\end{equation}

Keep $t,s$ fixed and set
\begin{equation}
 \widehat\alpha:=\frac{tL}{1-t},\qquad
 \widehat\beta:=s-\widehat\alpha.
 \label{eq:o3-normalized-constants}
\end{equation}
Then $0<\widehat\alpha<\alpha$, $\widehat\beta>\beta>0$, and
$\widehat\alpha+\widehat\beta=s<z$.  Changing these affine constants is a
translation and preserves $V_0\in\interior(\widehat T)$.  Its third
inequality remains \eqref{eq:o3-wedge-cut}; $\widehat U=0$ only at the
endpoint $(L/(1-t),0)$ of that common boundary and $\widehat V>0$.  Therefore
\[
 \widehat T\cap W=T\cap W,
\]
and their interiors have the same trace on $W$.  The local wedge calculation
identifies these with the closed and open traces in $H$.

Equation~\eqref{eq:o3-normalized-constants} places $Q_1$ on the boundary of
$H$; $Q_0$ remains outside and the second strict inequality in
\eqref{eq:o3-redundancy-conditions} keeps $Q_2$ outside.  Hence
$(o(\widehat T),n(\widehat T))=(2,0)$.  Since the three reach segments lie in
$H$, their open and closed traces, and therefore $A_i,B_i,C_i$ and
$A_i+B_i>1$, are unchanged.
\end{proof}

\subsection{The local admissible set and radial envelope}
\label{sec:admissible-set-derivation}

Let $u,v$ be unit vectors meeting at $120^\circ$ and define
\[
 K(a,b,c):=\conv\{0,au,bv,c(u+v)\}.
\]
This is the translate of the body normalization at $V_0$, with
$u=V_5-V_0$ and $v=V_1-V_0$.
The exact local admissible set is
\begin{equation}
 \mathcal A:=
 \{(a,b,c)\in[0,1]^3:\Lambda(K(a,b,c))\le1\}.
 \label{eq:universal-admissible-gauge}
\end{equation}
This set records the three reach lower bounds that a closed unit V triangle
can realize.

\begin{proposition}[Exact local admissible set]
\label{prop:new-exact-local-set}
\label{prop:exact-admissible-set}
Put
\[
 s:=a+b,\qquad d:=a^2+ab+b^2,\qquad
 m:=\min\{a,b\},\qquad M:=\max\{a,b\},
\]
\[
 q:=s^4-s^2+ab,
\]
and
\begin{align*}
 F_L&:=c^4-c^2+mc-m^2,\\
 F_T&:=(s^2-1)c^2+Mc-M^2,\\
 F_S&:=(m^2-1)c^2+(2mM^2+M)c+M^4-M^2.
\end{align*}
Then
\begin{equation}
 \mathcal A=\mathcal A_L\cup\mathcal A_T\cup\mathcal A_S,
 \label{eq:derived-admissible-cells}
\end{equation}
where all variables lie in $[0,1]$ and
\begin{align*}
 \mathcal A_L&:=\{d\le1,\ s\le1,\ q\le0,\ F_L\le0\},\\
 \mathcal A_T&:=\{d\le1,\ s\le1,\ q\ge0,\ F_T\le0,\ c\le2M\},\\
 \mathcal A_S&:=\{d\le1,\ s>1,\ F_S\le0,\ c\le1/2\}.
\end{align*}
The set is coordinatewise down-closed.  Define its maximal radial demand by
\[
 c_{\max}(a,b):=\max\{c\in[0,1]:(a,b,c)\in\mathcal A\}.
\]
Then
\begin{equation}
 c_{\max}(a,b)=
 \begin{cases}
 1,&a=b=0,\\
 C_L(m),&s\le1,\ q\le0,\ (a,b)\ne(0,0),\\
 \dfrac{2M}{1+\sqrt{4s^2-3}},&s\le1,\ q>0,\\[3mm]
 \dfrac{M(2mM+1)-M\sqrt{4d-3}}{2(1-m^2)},&s>1,
 \end{cases}
 \label{eq:radial-envelope}
\end{equation}
where $C_L(m)$ is the unique root in $[\sqrt3/2,1]$ of
\[
 z^4-z^2+mz-m^2=0.
\]
\end{proposition}

\begin{proof}
First assume $s>0$ and $cs\le ab$.  When $a,b>0$,
\[
 c(u+v)=\frac ca(au)+\frac cb(bv)
       +\left(1-\frac ca-\frac cb\right)0
 \in\conv\{0,au,bv\};
\]
the zero-coordinate cases follow by continuity.  The distance from $au$ to
$bv$ is $\sqrt d$.  With $T=-bu-av$, the triangle
$\conv\{au,bv,T\}$ is equilateral of side $\sqrt d$ and
\[
 0=\frac{b^2}{d}(au)+\frac{a^2}{d}(bv)+\frac{ab}{d}T.
\]
Thus the least enclosure side is $\sqrt d$.  When $a=b=0$, it is $c$.

Now assume $cs>ab$.  The hull has cyclic vertices $0,bv,c(u+v),au$.
Proposition~\ref{prop:universal-enclosure-gauge} reduces its least enclosure
side to its four edge normals.  For the edges through $0$ one obtains
\[
 L_{0A}=b+\max\{a,c\},
 \qquad L_{B0}=a+\max\{b,c\}.
\]
For the edge from $au$ to $c(u+v)$ put
\[
 R_a:=\sqrt{a^2-ac+c^2},
 \qquad n_0:=\frac1{2R_a}(\sqrt3a,2c-a).
\]
The three support values in the rotated directions are
\[
 \frac{\sqrt3ac}{2R_a},\qquad
 \frac{\sqrt3a(a-c)_+}{2R_a},\qquad
 \frac{\sqrt3c\max\{b,(c-a)_+\}}{2R_a}.
\]
Consequently
\[
 L_{AC}=
 \begin{cases}
 (a^2+bc)/R_a,&a\ge c,\\
 c(a+b)/R_a,&a<c\le a+b,\\
 c^2/R_a,&c>a+b,
 \end{cases}
\]
and $L_{CB}$ is obtained by interchanging $a,b$.  Hence
\[
 \Lambda(K(a,b,c))
 =\min\{L_{0A},L_{B0},L_{AC},L_{CB}\}.
\]

Assume by symmetry that $a=m\le b=M$.  Reading the active support points in
these four formulas gives
\[
\begin{array}{cc|ccc|c}
\multicolumn{2}{c|}{\text{range}}
 &\multicolumn{3}{c|}{\text{active supports in cyclic order}}
 &\text{unit frontier}\\
s&q&\text{first}&\text{second}&\text{third}&\text{equation}\\
\hline
s<1&q<0&\{0\}&\{c(u+v)\}&\{au,c(u+v)\}&F_L=0\\
s<1&q>0&\{0\}&\{bv,c(u+v)\}&\{au\}&F_T=0\\
\multicolumn{2}{c|}{s=1}&\{0,bv\}&\{bv,c(u+v)\}&\{au\}&c=M\\
\multicolumn{2}{c|}{s>1}&\{bv\}&\{bv,c(u+v)\}&\{au\}&F_S=0
\end{array}
\]
The inactive-support inequalities reduce to $q\le0$, $q\ge0$, and $s>1$.
The positive unit-frontier equations square to $F_L=0,F_T=0,F_S=0$.

It remains to select the geometric components introduced by squaring.  In
Cell $T$ the roots are
\[
 c_-:=\frac{2M}{1+\sqrt{4s^2-3}},
 \qquad
 c_+:=\frac{2M}{1-\sqrt{4s^2-3}}.
\]
The polynomial inequality also has a formal high component, while the
geometric fiber is the component joined to $c=0$.  Since
\[
 c_-\le2M\le c_+,
\]
the exact selector is $c\le2M$.

In Cell $S$,
\[
 F_S(a,b,0)=M^2(M^2-1)<0,
 \qquad F_S(a,b,M)=M^2(s^2-1)>0.
\]
Moreover $4F_S(a,b,1/2)$ increases with $m$, and at $m=1-M$ equals
$M^2(2M-1)^2$.  Because $s>1$, one has $F_S(a,b,1/2)>0$.  The quadratic
opens downward, so $c\le1/2$ selects exactly the smaller-root component
containing $0$.

In Cell $L$, the fiber from $c=0$ reaches one root in
$[\sqrt3/2,1]$; uniqueness there follows from
$4c^3-2c+m>0$.  At $q=0$ the $L$ and $T$ formulas meet at $c=s$.
This proves \eqref{eq:derived-admissible-cells}, including all equality and
degenerate cases by continuity.  Convexity of any containing triangle proves
coordinatewise down-closedness.  Solving the three selected frontier
equations gives \eqref{eq:radial-envelope}.
\end{proof}

For later variational use, the raw forward envelope retains its established
notation
\[
 M_c(a):=\max\{b\in[0,1]:(a,b,c)\in\mathcal A\}.
\]
No closed-form evaluation of $M_c$ is needed in the geometric body.

\subsection{Midpoint caliper optimization}

\begin{lemma}[Four-point self-midpoint caliper]
\label{lem:app-self-midpoint-caliper}
If $a+b>1$, then
\[
 \Lambda\!\left(\conv\left\{0,au,\frac{u+v}{2},bv\right\}\right)>1.
\]
This supplies Lemma~\ref{lem:self-midpoint}.
\end{lemma}

\begin{proof}
Put
\[
 K:=\conv\left\{0,au,\frac{u+v}{2},bv\right\}.
\]
By Theorem~\ref{thm:cert-caliper}, it suffices to check a normal to each of
the four hull edges.  For the edge from $0$ to $au$, the sum of the three
supports is
\[
 S_{0A}=\max\left\{\frac{\sqrt3}{4},\frac{\sqrt3a}{2}\right\}
          +\frac{\sqrt3b}{2}>\frac{\sqrt3}{2}.
\]
If $a\ge1/2$, this is $(\sqrt3/2)(a+b)$; if $a<1/2$, then $b>1/2$.
The edge through $bv$ is symmetric.

For the edge from $au$ to $(u+v)/2$, put
$D_a:=\sqrt{4a^2-2a+1}$.  Projection gives
\[
 S_{AM}=
 \begin{cases}
 \dfrac{\sqrt3(a+b)}{2D_a},&a\le1/2,\\[4pt]
 \dfrac{\sqrt3(2a^2+b)}{2D_a},&a>1/2.
 \end{cases}
\]
The first value exceeds $\sqrt3/2$ because $D_a\le1$.  In the second case,
$b>1-a$ and
\[
 (2a^2-a+1)^2-D_a^2=a^2(2a-1)^2\ge0,
\]
so it again exceeds $\sqrt3/2$.  The remaining edge is symmetric.  Formula
\eqref{eq:universal-lambda} now gives $\Lambda(K)>1$.

If $a+b=1$ while all four points are contained with positive margin, both
boundary reach demands can be increased slightly, contradicting the strict
case just proved.
\end{proof}

\subsection{Signed-center optimization}
\label{sec:signed-ce12-reductions}

The CE1 and CE2 cases share one signed cell.  The parameters below are used
only in the appendices; the geometric body uses the exact C triangle type and
midpoint interfaces.

\begin{proposition}[Signed CE1/CE2 center normal form]
\label{prop:signed-center-normal-form}
Let $T_C$ be a closed unit equilateral triangle with
$O\in\interior(T_C)$ and a positive-length boundary trace.  Normalize such a
trace to $e_{0,1}$ and use affine coordinates
\[
 X=V_0+b(V_1-V_0)+a(V_5-V_0).
\]
There are parameters
\[
 \begin{gathered}
  0<R<1,\qquad W:=1-R,\qquad E:=\sqrt{1-RW},\\
  \eta:=1-E,\qquad P:=E(1-E),
 \end{gathered}
\]
and nonnegative slacks $\alpha,\delta$, with
$k:=\eta+\alpha+\delta$, such that
\begin{equation}
\begin{aligned}
 F_0&=R+\alpha-a+Wb,\\
 F_1&=Rb+Wa-k,\\
 F_2&=W+\delta-b+Ra,
\end{aligned}
\qquad T_C=\{F_0,F_1,F_2\ge0\}.
\label{eq:signed-side-slacks}
\end{equation}
Define
\begin{equation}
 \Delta_R:=P-\alpha-W\delta,
 \qquad \Delta_L:=P-R\alpha-\delta.
 \label{eq:signed-surpluses}
\end{equation}
Then
\begin{equation}
 T_C\cap e_{0,1}=\left[\frac{k}{R},W+\delta\right],
 \qquad
 \Haus(T_C\cap e_{0,1})=\frac{\Delta_R}{R},
 \label{eq:signed-right-trace}
\end{equation}
and the possible companion trace is
\begin{equation}
 T_C\cap e_{5,0}=\left[\frac{k}{W},R+\alpha\right]
 \quad\text{when }\Delta_L>0,
 \qquad
 \Haus(T_C\cap e_{5,0})=\frac{[\Delta_L]_+}{W}.
 \label{eq:signed-left-trace}
\end{equation}
The normalized trace requires $\Delta_R>0$, and
\begin{equation}
 T_C\text{ is }\CEone\iff\Delta_L\le0,
 \qquad
 T_C\text{ is }\CEtwo\iff\Delta_L>0.
 \label{eq:signed-classification}
\end{equation}
The six radial exits from $O$ are
\begin{equation}
\begin{aligned}
 d_0^C&=E-\alpha-\delta,&
 d_1^C&=\frac{\delta}{R},&
 d_2^C&=\delta,\\
 d_3^C&=\min\left\{\frac{\alpha}{R},\frac{\delta}{W}\right\},&
 d_4^C&=\alpha,&
 d_5^C&=\frac{\alpha}{W}.
\end{aligned}
\label{eq:signed-center-exits}
\end{equation}
Moreover,
\begin{equation}
 T_C\cap\{M_0,\ldots,M_5\}=\{M_0\}.
 \label{eq:signed-unique-center-midpoint}
\end{equation}
For the closure of an original open C triangle, $\alpha,\delta>0$.
\end{proposition}

\begin{proof}
Write $T_C\cap e_{0,1}=[s,t]$ with $0\le s<t\le1$.  The two sides through
the endpoints have a positive slope ratio.  Reflect across the perpendicular
bisector if necessary, call the ratio $R\le1$, and put $W=1-R$.  The three
side equations are
\[
 F_1=Rb+Wa-Rs,\qquad
 F_2=-b+Ra+t,\qquad
 F_0=Wb-a+E+Rs-t,
\]
where $E=\sqrt{1-RW}$ is the unit-side relation.  If $R=1$, then
$F_0(O)=s-t<0$, so $0<R<1$.

Set $\alpha:=F_0(O)$ and $\delta:=F_2(O)$.  Then
\[
 t=W+\delta,\qquad Rs=\eta+\alpha+\delta=k,
\]
which gives \eqref{eq:signed-side-slacks}.  The gradients are the three
equilateral side-normal directions and $F_0+F_1+F_2=E$.

On $e_{0,1}$, $a=0$ and the active inequalities are
$Rb\ge k$ and $b\le W+\delta$.  Thus its length is
\[
 W+\delta-\frac{k}{R}
 =\frac{P-\alpha-W\delta}{R}.
\]
On $e_{5,0}$, $b=0$ and the active inequalities are
$Wa\ge k$ and $a\le R+\alpha$, giving signed length
\[
 R+\alpha-\frac{k}{W}
 =\frac{P-R\alpha-\delta}{W}.
\]
This proves the two trace formulas.

The positive right trace implies
\[
 \delta<\frac{P}{W}=\frac{ER}{1+E}<\frac R2.
\]
On $e_{1,2}$, parametrized by $b=1+q,a=q$, one has
$F_2=\delta-R-Wq<0$.  The center classification leaves $e_{5,0}$ as the only
possible companion edge, and its signed length proves
\eqref{eq:signed-classification}; $\Delta_L=0$ is only point contact.

Substituting the six ray parametrizations into
\eqref{eq:signed-side-slacks} gives \eqref{eq:signed-center-exits}.  Finally,
\[
 W(\alpha+\delta)\le\alpha+W\delta<P
\]
implies
\[
 E-\alpha-\delta>E-\frac PW
 =\frac{E(E-R)}W>\frac12,
\]
where the last inequality is equivalent to $(E-R)^2>0$.  Also
\[
 \alpha<P<\min\left\{\frac R2,\frac W2\right\},
 \qquad\delta<\frac R2.
\]
Therefore $d_0^C>1/2$ and $d_i^C<1/2$ for $i=1,\ldots,5$, proving the
midpoint formula.  Interior containment of $O$ in the original open triangle
makes all three side slacks positive, so $\alpha,\delta>0$.
\end{proof}

\begin{lemma}[CE2 total slack]
\label{lem:new-signed-ce2-total-slack}
In the CE2 sign range, with $T:=\alpha+\delta$,
\begin{equation}
 T<\frac{\eta}{E}.
 \label{eq:signed-total-slack}
\end{equation}
\end{lemma}

\begin{proof}
The CE2 inequalities are
$\alpha+W\delta<P$ and $R\alpha+\delta<P$.  Multiply them by $W$ and $R$
and add.  Since $W+R^2=W^2+R=E^2$, this gives
$E^2T<P=E\eta$.
\end{proof}

\begin{corollary}[Common center-boundary formula]
\label{cor:signed-center-boundary}
The full center-boundary contribution is
\begin{equation}
 L_{\partial H}(T_C)
 =\frac{\Delta_R}{R}+\frac{[\Delta_L]_+}{W}.
 \label{eq:signed-boundary-length}
\end{equation}
In CE1,
\[
 L_{\partial H}(T_C)\le P\le\frac{\sqrt3}{2}-\frac34,
\]
and in CE2,
\[
 L_{\partial H}(T_C)
 =\frac{E}{1+E}-\frac{E^2(\alpha+\delta)}{RW}<\frac12.
\]
\end{corollary}

\begin{proof}
The first formula sums \eqref{eq:signed-right-trace} and
\eqref{eq:signed-left-trace}.  In CE1,
$R\alpha+\delta\ge P$; multiplication by $W$, together with
$\alpha\ge RW\alpha$, gives $\alpha+W\delta\ge WP$.  Thus
$\Delta_R\le RP$ and the contribution is at most $P$.  Since
$E\ge\sqrt3/2$ and $E(1-E)$ decreases there, the stated cap follows.

In CE2, direct addition gives
\[
 L_{\partial H}(T_C)
 =\frac{P-E^2(\alpha+\delta)}{RW}
 =\frac{E}{1+E}-\frac{E^2(\alpha+\delta)}{RW}<\frac12.
\]
\end{proof}

\begin{lemma}[Adjacent-edge diameter transfer]
\label{lem:signed-diameter-transfer}
For $0\le q\le1$, the zero-radial forward envelope is
\[
 M_0(q)=\frac{-q+\sqrt{4-3q^2}}2.
\]
Here $M_0$ without an argument remains the radial midpoint.  If a unit
equilateral triangle reaches parameters $q,b$ on adjacent boundary edges,
then $b\le M_0(q)$.  The function is strictly decreasing and, for $q>0$,
\[
 M_0(q)<1-\frac q2,
 \qquad 1-M_0(q)>\frac q2.
\]
\end{lemma}

\begin{proof}
The squared distance between the boundary points is $q^2+qb+b^2\le1$.
Solving the quadratic in $b$ gives the formula.  Differentiation gives
\[
 M_0'(q)=\frac{-1-3q/\sqrt{4-3q^2}}2<0.
\]
For $q>0$, $\sqrt{4-3q^2}<2$, which gives both strict comparisons.
\end{proof}
